\documentclass[11pt,a4paper]{article}
% Packages
\usepackage[utf8]{inputenc}
\usepackage[T1]{fontenc}
\usepackage[margin=1in]{geometry}
\usepackage{listings}
\usepackage{xcolor}
\usepackage{booktabs}
\usepackage{hyperref}
\usepackage{fancyhdr}
\usepackage{tcolorbox}
\usepackage{enumitem}
\usepackage{titlesec}
% Colors
\definecolor{codegreen}{rgb}{0,0.6,0}
\definecolor{codegray}{rgb}{0.5,0.5,0.5}
\definecolor{codepurple}{rgb}{0.58,0,0.82}
\definecolor{backcolour}{rgb}{0.95,0.95,0.92}
\definecolor{cmdcolor}{rgb}{0.2,0.2,0.6}
\definecolor{warningcolor}{rgb}{0.8,0.2,0.2}
\definecolor{tipcolor}{rgb}{0.1,0.5,0.1}
% Code listing style
\lstdefinestyle{bashstyle}{
    backgroundcolor=\color{backcolour},
    commentstyle=\color{codegreen},
    keywordstyle=\color{cmdcolor}\bfseries,
    numberstyle=\tiny\color{codegray},
    stringstyle=\color{codepurple},
    basicstyle=\ttfamily\small,
    breakatwhitespace=false,
    breaklines=true,
    captionpos=b,
    keepspaces=true,
    numbers=none,
    numbersep=5pt,
    showspaces=false,
    showstringspaces=false,
    showtabs=false,
    tabsize=2,
    frame=single,
    rulecolor=\color{codegray},
    xleftmargin=0.5em,
    xrightmargin=0.5em,
}
\lstset{style=bashstyle}
% Custom boxes
\newtcolorbox{warningbox}{
    colback=red!5!white,
    colframe=red!75!black,
    title=\textbf{Warning},
    fonttitle=\bfseries
}
\newtcolorbox{tipbox}{
    colback=green!5!white,
    colframe=green!50!black,
    title=\textbf{Tip},
    fonttitle=\bfseries
}
\newtcolorbox{notebox}{
    colback=blue!5!white,
    colframe=blue!50!black,
    title=\textbf{Note},
    fonttitle=\bfseries
}
% Header/Footer
\pagestyle{fancy}
\fancyhf{}
\rhead{HPC Notes: The Mill}
\lhead{Missouri S\&T}
\rfoot{Page \thepage}
% Title formatting
\titleformat{\section}{\Large\bfseries\color{cmdcolor}}{\thesection}{1em}{}
\titleformat{\subsection}{\large\bfseries}{\thesubsection}{1em}{}
% Document info
\title{\textbf{HPC Notes: The Mill Cluster}\\[0.5em]\large Missouri University of Science \& Technology}
\author{Research Computing Reference Guide}
\date{Last Updated: \today}
\begin{document}
\maketitle
\tableofcontents
\newpage
%=============================================================================
\section{Connection Info}
%=============================================================================
\begin{lstlisting}[language=bash]
# SSH connection
ssh mg6f4@mill-login-p2.itrss.mst.edu
# Or use the short alias (if configured in ~/.ssh/config)
ssh mill
\end{lstlisting}
\begin{tipbox}
Add this to your local \texttt{\textasciitilde/.ssh/config} for easier access:
\begin{lstlisting}[language=bash]
Host mill
    HostName mill-login-p2.itrss.mst.edu
    User mg6f4
\end{lstlisting}
\end{tipbox}
%=============================================================================
\section{Quick Reference}
%=============================================================================
\begin{lstlisting}[language=bash]
# Check your jobs
squeue -u $USER
# Cancel a job
scancel <JOBID>
# Check job history
sacct -u $USER --starttime=$(date -d "3 days ago" +%Y-%m-%d) \
      --format=JobID,JobName%15,State,Elapsed,ExitCode
# Check cluster status
sinfo -o "%5D %4c %8m %28f %35G"
\end{lstlisting}
%=============================================================================
\section{Module Loading (deal.II)}
%=============================================================================
\begin{lstlisting}[language=bash]
# Load the GCC-based deal.II module (REQUIRED: must use this one)
module load dealii/9.7.1_gcc_12.3.0_mvapich
# Verify
module list
\end{lstlisting}
\begin{warningbox}
Do NOT use \texttt{dealii/9.7.1\_umfpack\_mumps} or other clang-based modules.
The deal.II library was built with GCC, but those modules load clang as the
compiler, causing an \textbf{ABI mismatch} (different C++ name mangling for
\texttt{std::is\_same\_v} and similar templates). Use the GCC module above.
\end{warningbox}
\begin{notebox}
The deal.II installation on The Mill is \textbf{2D only}. If your code has
3D template instantiations (e.g., \texttt{template class MyClass<3>}),
comment them out before building.

Additionally, \texttt{VectorTools::project} is not instantiated for
\texttt{TrilinosWrappers::MPI::Vector}. Use
\texttt{VectorTools::interpolate} followed by
\texttt{constraints.distribute()} instead.
\end{notebox}

\subsection{Available Libraries and Packages}
The \texttt{dealii/9.7.1\_gcc\_12.3.0\_mvapich} module includes:

\begin{table}[h]
\centering
\begin{tabular}{@{}lll@{}}
\toprule
\textbf{Library} & \textbf{Status} & \textbf{Notes} \\
\midrule
Trilinos & \checkmark & Core framework \\
\quad Amesos2 & \checkmark & Direct solver interface (KLU2, Basker) \\
\quad Belos & \checkmark & Krylov iterative solvers \\
\quad IFPACK2 & \checkmark & ILU preconditioners \\
\quad ML/MueLu & \texttimes & AMG not available \\
PETSc & \checkmark & With Hypre, SuperLU\_dist \\
\quad SuperLU\_dist & \checkmark & Parallel direct solver (via PETSc) \\
p4est & \checkmark & Parallel AMR \\
METIS & \checkmark & Graph partitioning \\
HDF5 & \checkmark & Data I/O \\
Kokkos & \checkmark & Performance portability \\
Ginkgo & \checkmark & GPU-capable solvers \\
MUMPS & \texttimes & Not available \\
UMFPACK & \texttimes & Not available \\
\bottomrule
\end{tabular}
\caption{Library availability in the GCC deal.II module}
\end{table}

\begin{warningbox}
The only Trilinos direct solver available is \textbf{KLU} (serial, slow).
The old Amesos interface (\texttt{Amesos\_Superlu}, \texttt{Amesos\_Mumps})
will fail. Use Amesos2 or switch to PETSc with SuperLU\_dist for
parallel direct solves.
\end{warningbox}

\begin{notebox}
To check what is enabled in the deal.II build:
\begin{lstlisting}[language=bash]
grep -i "DEAL_II_WITH_" \
  /share/apps/src/dealii/dealii-9.7.1/dealii-9.7.1_install/\
  include/deal.II/base/config.h | grep -v "undef"
\end{lstlisting}
To check PETSc packages:
\begin{lstlisting}[language=bash]
grep -i "PETSC_HAVE_" \
  /share/apps/common/SPACK_21/spack/opt/spack/\
  linux-almalinux8-x86_64/gcc-12.3.0/petsc-*/\
  include/petscconf.h
\end{lstlisting}
\end{notebox}
\subsection{Available deal.II Modules}
\begin{lstlisting}[language=bash]
module avail dealii
# dealii/9.5.2_gcc_12.3.0_mvapich
# dealii/9.7.1_gcc_12.3.0_mvapich   <-- USE THIS ONE
# dealii/9.7.1_umfpack_mumps         (clang ABI mismatch!)
# dealii/9.7.1_umfpack_openmpi
# dealii/9.7.1_umfpack               (multi-node MPI, needs source)
# dealii/9.7.1_umfpack2
\end{lstlisting}

\begin{notebox}
The \texttt{dealii/9.7.1\_umfpack} module supports multi-node MPI and has UMFPACK,
but requires sourcing the spack build environment after loading:
\begin{lstlisting}[language=bash]
module load dealii/9.7.1_umfpack
source /share/apps/common/SPACK_v110_release_a2/spack/opt/\
spack/linux-x86_64/dealii-9.7.1-kr766e7zm3gq6q7t67vefndkrhcyvmrw/\
.spack/spack-build-env.txt
\end{lstlisting}
Check for ABI compatibility before using this module with your code.
The \texttt{dealii/9.7.1\_umfpack\_mumps} module has MUMPS and UMFPACK but
was built with clang, causing link-time ABI mismatch with GCC-compiled code.
A rebuild with GCC has been requested.
\end{notebox}
\subsection{Other Module Commands}
\begin{lstlisting}[language=bash]
module avail              # List available modules
module list               # List loaded modules
module purge              # Unload all modules
module unload <name>      # Unload specific module
module swap <old> <new>   # Swap one module for another
\end{lstlisting}
%=============================================================================
\section{Partitions and Time Limits}
%=============================================================================
\begin{table}[h]
\centering
\begin{tabular}{@{}llll@{}}
\toprule
\textbf{Partition} & \textbf{Max Time} & \textbf{Access} & \textbf{Notes} \\
\midrule
requeue & 2 days & Open & May be preempted \\
general & 2 days & Open & Standard jobs \\
gpu & 2 days & Open & GPU jobs \\
interactive & 4 hours & Open & Interactive work \\
\textbf{hpcc} & \textbf{28 days} & \textbf{Leased} & \textbf{No 2-day limit!} \\
priority & 28 days & Leased & Priority access \\
\bottomrule
\end{tabular}
\caption{Partition time limits}
\end{table}
\begin{tipbox}
To bypass the 2-day time limit on \texttt{general}, use
\texttt{--partition=hpcc} (requires leased access).
The \texttt{hpcc} partition has \textbf{no default time limit}
--- you set \texttt{--time} yourself (up to 28 days).
\begin{lstlisting}[language=bash]
#SBATCH --partition=hpcc
#SBATCH --time=4-00:00:00    # 4 days (you choose)
\end{lstlisting}
\end{tipbox}
%=============================================================================
\section{Interactive Sessions}
%=============================================================================
\begin{warningbox}
Never run computations on the login node! Always use interactive sessions or batch jobs.
\end{warningbox}
\subsection{Start Interactive Session}
\begin{lstlisting}[language=bash]
# Standard (hpcc partition, 1 hour, 8GB, 4 cores)
salloc --partition=hpcc --time=1:00:00 --mem=8G --ntasks=4
# For compiling (more cores for parallel make)
salloc --partition=hpcc --time=1:00:00 --mem=16G --ntasks=16
# With X11 forwarding (for GUI apps)
salloc --partition=hpcc --time=1:00:00 --mem=8G --ntasks=4 --x11
\end{lstlisting}
\subsection{For MPI Jobs in Interactive Session}
\begin{lstlisting}[language=bash]
export SLURM_OVERLAP=1
mpirun -np 4 ./your_executable
\end{lstlisting}
\subsection{Exit Interactive Session}
\begin{lstlisting}[language=bash]
exit
\end{lstlisting}
%=============================================================================
\section{Compiling on Compute Nodes}
%=============================================================================
\begin{warningbox}
Always compile on compute nodes to avoid ``Illegal instruction'' errors!
The login node has different CPU architecture than compute nodes.
\end{warningbox}
\begin{lstlisting}[language=bash]
# 1. Get interactive session
salloc --partition=hpcc --time=1:00:00 --mem=16G --ntasks=16
# 2. Load module
module load dealii/9.7.1_gcc_12.3.0_mvapich
# 3. Navigate to project and create build directory
cd ~/Decoupled-HPC/Decoupled
mkdir -p build && cd build
# 4. Configure and build
cmake -DCMAKE_BUILD_TYPE=Release ..
make -j16
# 5. Exit when done
exit
\end{lstlisting}
\begin{notebox}
Project locations on the cluster (sparse git checkouts):
\begin{itemize}[nosep]
    \item Decoupled: \texttt{\textasciitilde/Decoupled-HPC/Decoupled/}
    \item Semi\_Coupled: \texttt{\textasciitilde/Semi\_Coupled-HPC/Semi\_Coupled/}
\end{itemize}
To update: \texttt{cd \textasciitilde/<project>-HPC/<project> \&\& git pull}
\end{notebox}
%=============================================================================
\section{Storage Locations}
%=============================================================================
\begin{table}[h]
\centering
\begin{tabular}{@{}llll@{}}
\toprule
\textbf{Location} & \textbf{Path} & \textbf{Size} & \textbf{Purpose} \\
\midrule
Home & \texttt{/home/mg6f4} & 50 GB & Code, scripts \\
Scratch & \texttt{/share/ceph/scratch/mg6f4} & Large & Run output, temp files \\
Local scratch & \texttt{/local/scratch} & 1.5 TB & Node-local temp \\
\bottomrule
\end{tabular}
\caption{Storage locations on The Mill}
\end{table}
\begin{warningbox}
Nothing is backed up! Always keep copies of important data elsewhere.
Scratch may be cleaned during maintenance.
\end{warningbox}
\begin{lstlisting}[language=bash]
# Go to scratch
cd /share/ceph/scratch/$USER
# or use the shortcut
cdsc
\end{lstlisting}
%=============================================================================
\section{Simulation Output Locations}
%=============================================================================
The ferrofluid solver writes output to several locations depending on the run type:
\subsection{Decoupled Project Directory Structure}
\begin{lstlisting}[language=bash]
~/Decoupled-HPC/Decoupled/
|-- build/                    # Executables
|   |-- bin0_hedgehog         # BIN-0: original Kelvin (H=grad(phi))
|   |-- bin1_hedgehog         # BIN-1: fixed Kelvin (H=h_a+grad(phi))
|   `-- drivers/
|       `-- ferrofluid_decoupled  # Main executable
|-- hpc/                      # HPC submission scripts
|   |-- hedgehog_array.sub    # SLURM job array (12 tests)
|   |-- hedgehog_tests.dat    # Test configuration matrix
|   `-- build_binaries.sh     # Build automation script
|-- logs/                     # SLURM job output logs
|   |-- T01_<jobid>.out       # Test 1 stdout
|   |-- T01_<jobid>.err       # Test 1 stderr
|   `-- ...
`-- Results/                  # Simulation output (PVTU, CSV)
    `-- <timestamp>_<config>/
        |-- phase_field_*.pvtu
        |-- velocity_*.pvtu
        |-- magnetization_*.pvtu
        |-- magnetic_potential_*.pvtu
        |-- energy.csv
        |-- convergence.csv
        |-- diagnostics.csv
        `-- warnings.csv
\end{lstlisting}
\subsection{Semi\_Coupled Project Directory Structure}
\begin{lstlisting}[language=bash]
~/Semi_Coupled-HPC/Semi_Coupled/
|-- cmake-build-release/      # Executables
|   `-- ferrofluid             # Main executable
|-- hpcc_run/                  # HPC submission scripts
|   |-- submit.sub             # Generic SLURM array script
|   |-- rosensweig_debug.dat   # Debug test matrix (13 tests)
|   `-- rosensweig_tests.dat   # Original test matrix
|-- logs/                      # SLURM job output logs
|   |-- rosen_T01_<jobid>.out
|   |-- rosen_T01_<jobid>.err
|   `-- ...
`-- Results/                   # Simulation output
    `-- <timestamp>_<run_name>/
        |-- solution_*.pvtu
        |-- energy.csv
        |-- diagnostics.csv
        |-- magnetic.csv
        |-- flow.csv
        |-- forces.csv
        |-- interface.csv
        |-- phase_field.csv
        |-- solver.csv
        |-- timing.csv
        `-- run_info.txt
\end{lstlisting}

\begin{notebox}
The Semi\_Coupled generic submission script takes the .dat file as an argument:
\begin{lstlisting}[language=bash]
cd ~/Semi_Coupled-HPC/Semi_Coupled
mkdir -p logs
sbatch --array=1-13 hpcc_run/submit.sub hpcc_run/rosensweig_debug.dat
\end{lstlisting}
\end{notebox}

\subsection{Output File Types}
\begin{table}[h]
\centering
\begin{tabular}{@{}lll@{}}
\toprule
\textbf{Extension} & \textbf{Description} & \textbf{Used For} \\
\midrule
\texttt{.pvtu} & Parallel VTK (master file) & ParaView visualization \\
\texttt{.vtu} & VTK per-process data & Referenced by .pvtu \\
\texttt{.csv} & Comma-separated values & Diagnostics, convergence rates \\
\texttt{.out} & SLURM stdout & Job logs, console output \\
\texttt{.err} & SLURM stderr & Error messages \\
\bottomrule
\end{tabular}
\caption{Common output file types}
\end{table}
\subsection{SLURM Log File Locations}
Job output files are written to the working directory specified in the submission script:
\begin{lstlisting}[language=bash]
# Check where a job writes output
scontrol show job <JOBID> | grep -E "WorkDir|StdOut|StdErr"
# List output files (most recent first)
ls -lt logs/*.out logs/*.err | head -10
\end{lstlisting}
\subsection{Finding Output Files}
\begin{lstlisting}[language=bash]
# Find all PVTU files (parallel VTK)
find ~/Decoupled-HPC -name "*.pvtu" | head -20
# Find all CSV results
find ~/Decoupled-HPC -name "*.csv"
# Find recent output directories
ls -lt ~/Decoupled-HPC/Decoupled/Results/ | head -10
\end{lstlisting}
%=============================================================================
\section{Creating Submission Scripts (.sub files)}
%=============================================================================
\subsection{Basic Template (with hpcc partition)}
\begin{lstlisting}[language=bash]
#!/bin/bash
#SBATCH --job-name=my_job          # Job name (shows in queue)
#SBATCH --output=my_job_%j.out     # Standard output (%j = job ID)
#SBATCH --error=my_job_%j.err      # Standard error
#SBATCH --partition=hpcc           # hpcc = 28-day access
#SBATCH --time=4-00:00:00          # Time limit (D-HH:MM:SS)
#SBATCH --mem=128G                 # Memory
#SBATCH --nodes=1                  # Number of nodes
#SBATCH --ntasks=16                # Number of MPI tasks
#SBATCH --mail-type=BEGIN,END,FAIL # Email notifications
# Load modules
module load dealii/9.7.1_gcc_12.3.0_mvapich
# Change to working directory
cd $SLURM_SUBMIT_DIR
# Run executable
mpirun -np $SLURM_NTASKS ./build/drivers/ferrofluid_decoupled --flags
\end{lstlisting}
\begin{warningbox}
Without \texttt{--partition=hpcc}, jobs are limited to \textbf{2 days}
on the \texttt{general} partition. Add \texttt{--partition=hpcc} to
run up to \textbf{28 days}.
\end{warningbox}
\subsection{Hedgehog Test Array (12 parallel jobs)}
\begin{lstlisting}[language=bash]
#!/bin/bash
#SBATCH --job-name=hedgehog
#SBATCH --partition=hpcc
#SBATCH --array=1-12              # 12 tests in parallel
#SBATCH --nodes=1                 # Each test on its own node
#SBATCH --ntasks=16               # 16 MPI ranks per test
#SBATCH --time=4-00:00:00
#SBATCH --mem=128G
#SBATCH --output=logs/T%02a_%j.out
#SBATCH --error=logs/T%02a_%j.err
module load dealii/9.7.1_gcc_12.3.0_mvapich
cd $SLURM_SUBMIT_DIR
# Read test config for this array task
LINE=$(awk -v line=${SLURM_ARRAY_TASK_ID} \
    '!/^#/ && !/^$/ {n++; if (n==line) print}' hpc/hedgehog_tests.dat)
BINARY=$(echo "$LINE" | awk '{print $1}')
FLAGS=$(echo "$LINE" | cut -d' ' -f2-)
# Map binary name to path
if [ "$BINARY" = "bin0" ]; then EXE="./build/bin0_hedgehog"
elif [ "$BINARY" = "bin1" ]; then EXE="./build/bin1_hedgehog"
else echo "ERROR: Unknown binary"; exit 1; fi
mpirun -np $SLURM_NTASKS $EXE $FLAGS
\end{lstlisting}
\subsection{Common SBATCH Options}
\begin{table}[h]
\centering
\begin{tabular}{@{}lll@{}}
\toprule
\textbf{Option} & \textbf{Description} & \textbf{Example} \\
\midrule
\texttt{--job-name} & Job name & \texttt{--job-name=my\_sim} \\
\texttt{--time} & Wall time limit & \texttt{--time=4-00:00:00} \\
\texttt{--mem} & Total memory & \texttt{--mem=128G} \\
\texttt{--mem-per-cpu} & Memory per CPU & \texttt{--mem-per-cpu=4G} \\
\texttt{--ntasks} & Number of MPI tasks & \texttt{--ntasks=16} \\
\texttt{--nodes} & Number of nodes & \texttt{--nodes=1} \\
\texttt{--partition} & Queue/partition & \texttt{--partition=hpcc} \\
\texttt{--array} & Job array & \texttt{--array=1-12} \\
\texttt{--gres} & GPU resources & \texttt{--gres=gpu:2} \\
\bottomrule
\end{tabular}
\caption{Common SBATCH directives}
\end{table}
%=============================================================================
\section{Job Submission and Management}
%=============================================================================
\subsection{Submit Jobs}
\begin{lstlisting}[language=bash]
# Single job
sbatch my_job.sub
# Job array (e.g., hedgehog tests)
cd ~/Decoupled-HPC/Decoupled
mkdir -p logs
sbatch hpc/hedgehog_array.sub
# Multiple individual jobs
for f in *.sub; do sbatch "$f"; done
\end{lstlisting}
\subsection{Monitor Jobs}
\begin{lstlisting}[language=bash]
# Your jobs
squeue -u $USER
# Detailed format
squeue -u $USER --format="%.10i %.12j %.8T %.12M %.6D %R"
# All jobs (cluster status)
squeue
# Watch continuously (updates every 30 sec)
watch -n 30 'squeue -u $USER'
\end{lstlisting}
\subsection{Job History}
\begin{lstlisting}[language=bash]
# Recent jobs (last 3 days)
sacct -u $USER --starttime=$(date -d "3 days ago" +%Y-%m-%d) \
      --format=JobID,JobName%15,State,Elapsed,ExitCode
# Specific job details
sacct -j <JOBID> --format=JobID,JobName,State,Elapsed,MaxRSS,ExitCode
\end{lstlisting}
\subsection{Cancel Jobs}
\begin{lstlisting}[language=bash]
# Cancel single job
scancel <JOBID>
# Cancel all your jobs
scancel -u $USER
# Cancel by job name
scancel --name=hedgehog
# Cancel specific array tasks
scancel <JOBID>_[1-6]    # Cancel tasks 1-6 only
\end{lstlisting}
%=============================================================================
\section{Checking Job Output}
%=============================================================================
\begin{lstlisting}[language=bash]
# Find where job output goes
scontrol show job <JOBID> | grep -E "WorkDir|StdOut|StdErr"
# View hedgehog test logs
cat logs/T01_*.out     # Test 1 output
cat logs/T01_*.err     # Test 1 errors
# Tail a running job's output
tail -f logs/T01_*.out
# Check all test statuses at once
for i in $(seq -w 1 12); do
    echo "=== T$i ==="; tail -1 logs/T${i}_*.out 2>/dev/null
done
\end{lstlisting}
\subsection{Common Exit Codes}
\begin{table}[h]
\centering
\begin{tabular}{@{}ll@{}}
\toprule
\textbf{Code} & \textbf{Meaning} \\
\midrule
0 & Success \\
1 & General error \\
127 & Command not found (bad path or module) \\
137 & Killed (out of memory or time) \\
Illegal instruction & CPU mismatch (recompile on compute node!) \\
\bottomrule
\end{tabular}
\caption{Common exit codes and their meanings}
\end{table}
%=============================================================================
\section{File Transfer}
%=============================================================================
\subsection{SCP (from local machine)}
\begin{lstlisting}[language=bash]
# Upload
scp file.txt mg6f4@mill-login-p2.itrss.mst.edu:~/destination/
# Download
scp mg6f4@mill-login-p2.itrss.mst.edu:~/file.txt ./local/
# Recursive (directories)
scp -r folder/ mg6f4@mill-login-p2.itrss.mst.edu:~/destination/
\end{lstlisting}
\subsection{Rsync (better for large transfers)}
\begin{lstlisting}[language=bash]
# Upload with progress
rsync -avz --progress ./local_folder/ \
      mg6f4@mill-login-p2.itrss.mst.edu:~/remote_folder/
# Download
rsync -avz --progress \
      mg6f4@mill-login-p2.itrss.mst.edu:~/remote_folder/ ./local_folder/
\end{lstlisting}
\subsection{Download Hedgehog Results}
\begin{lstlisting}[language=bash]
# Run from LOCAL machine
cd /Users/mahdi/Projects/hpc_results
# Logs
rsync -avz --progress \
    mg6f4@mill-login-p2.itrss.mst.edu:~/Decoupled-HPC/Decoupled/logs/ \
    ./logs/
# Simulation results (PVTU + CSV)
rsync -avz --progress \
    mg6f4@mill-login-p2.itrss.mst.edu:~/Decoupled-HPC/Decoupled/Results/ \
    ./Results/
\end{lstlisting}
\subsection{Globus}
For large transfers, use Globus at \url{https://globus.org}
\begin{itemize}
    \item Collection name: ``Missouri S\&T Mill''
\end{itemize}
%=============================================================================
\section{Useful One-Liners}
%=============================================================================
\begin{lstlisting}[language=bash]
# Check disk usage in home
du -sh ~/*
# Check disk usage in scratch
du -sh /share/ceph/scratch/$USER/*
# Count PVTU files in Results
find ~/Decoupled-HPC/Decoupled/Results -name "*.pvtu" | wc -l
# Find recent output directories
ls -lt ~/Decoupled-HPC/Decoupled/Results/ | head -10
# Check if executables exist
ls -lh ~/Decoupled-HPC/Decoupled/build/bin*_hedgehog
# Check running job's resource usage
sstat -j <JOBID> --format=JobID,MaxRSS,MaxVMSize,AveCPU
# Find large files
find ~/Decoupled-HPC -size +100M -exec ls -lh {} \;
\end{lstlisting}
%=============================================================================
\section{Troubleshooting}
%=============================================================================
\subsection{``Illegal instruction (core dumped)''}
\begin{itemize}
    \item \textbf{Cause:} Compiled on login node, running on compute node with different CPU
    \item \textbf{Fix:} Recompile on a compute node (see Section 6)
\end{itemize}
\subsection{Undefined symbol / ABI mismatch}
\begin{itemize}
    \item \textbf{Cause:} Compiled with clang but deal.II was built with GCC
    \item \textbf{Symptom:} \texttt{undefined symbol: ...is\_same\_v...} during linking
    \item \textbf{Fix:} Use \texttt{module load dealii/9.7.1\_gcc\_12.3.0\_mvapich}
\end{itemize}
\subsection{3D template instantiation errors}
\begin{itemize}
    \item \textbf{Cause:} deal.II on The Mill is 2D only
    \item \textbf{Fix:} Comment out all \texttt{template class MyClass<3>} lines
\end{itemize}
\subsection{VectorTools::project linker error}
\begin{itemize}
    \item \textbf{Cause:} Not instantiated for \texttt{TrilinosWrappers::MPI::Vector}
    \item \textbf{Fix:} Replace with \texttt{VectorTools::interpolate} + \texttt{constraints.distribute()}
\end{itemize}
\subsection{Exit code 127}
\begin{itemize}
    \item \textbf{Cause:} Executable or command not found
    \item \textbf{Fix:} Check paths, ensure module is loaded in .sub file
\end{itemize}
\subsection{Job stuck in queue (PD status)}
\begin{itemize}
    \item \textbf{Cause:} Waiting for resources
    \item \textbf{Check:} \texttt{squeue -u \$USER} shows REASON column
\end{itemize}
\subsection{Out of memory (exit code 137)}
\begin{itemize}
    \item \textbf{Fix:} Increase \texttt{--mem} in submission script
\end{itemize}
\subsection{Job timeout}
\begin{itemize}
    \item \textbf{Fix:} Increase \texttt{--time} or use \texttt{--partition=hpcc} for up to 28 days
\end{itemize}

\subsection{Amesos\_Superlu / Amesos\_Mumps not found}
\begin{itemize}
    \item \textbf{Cause:} Trilinos was built without SuperLU or MUMPS. Only KLU is available.
    \item \textbf{Symptom:} \texttt{You tried to select the solver type <Amesos\_Superlu> but this solver is not supported}
    \item \textbf{Workaround:} The code falls back to \texttt{Amesos\_Klu} (serial, slow). For better performance, switch to PETSc with SuperLU\_dist or use Amesos2 interface.
\end{itemize}

\subsection{Jobs pending: ReqNodeNotAvail, Reserved for maintenance}
\begin{itemize}
    \item \textbf{Cause:} Job time limit extends past a scheduled maintenance window.
    \item \textbf{Check:} \texttt{scontrol show reservation} to see maintenance dates.
    \item \textbf{Fix:} Reduce \texttt{--time} so the job finishes before maintenance starts.
\end{itemize}

\subsection{Cannot extend running job time limit}
\begin{itemize}
    \item \textbf{Cause:} SLURM does not allow regular users to extend jobs past the originally submitted time limit via \texttt{scontrol update}.
    \item \textbf{Fix:} Submit with sufficient time from the start. If a job will hit its wall time, cancel and resubmit with a longer limit.
\end{itemize}

\subsection{Iterative solver crash (Unknown exception)}
\begin{itemize}
    \item \textbf{Cause:} AMG preconditioner (Trilinos ML/MueLu) not available in the current build.
    \item \textbf{Symptom:} Job crashes immediately after setup with ``Unknown exception''.
    \item \textbf{Fix:} Use direct solver instead of \texttt{--iterative}, or wait for MueLu to be added to the build.
\end{itemize}
%=============================================================================
\section{Quick Workflow Summary}
%=============================================================================
\begin{lstlisting}[language=bash]
# 1. Connect
ssh mill
# 2. Compile (do this on compute node!)
salloc --partition=hpcc --time=1:00:00 --mem=16G --ntasks=16
module load dealii/9.7.1_gcc_12.3.0_mvapich
cd ~/Decoupled-HPC/Decoupled
mkdir -p build && cd build
cmake -DCMAKE_BUILD_TYPE=Release ..
make -j16
exit
# 3. Submit (from login node)
cd ~/Decoupled-HPC/Decoupled
mkdir -p logs
sbatch hpc/hedgehog_array.sub
# 4. Monitor
squeue -u $USER
tail -f logs/T01_*.out
# 5. Check results
ls Results/
# 6. Download results (from LOCAL machine)
rsync -avz --progress \
    mg6f4@mill-login-p2.itrss.mst.edu:~/Decoupled-HPC/Decoupled/Results/ \
    ./local_results/
\end{lstlisting}
%=============================================================================
\section{Citation}
%=============================================================================
When publishing research that used The Mill, please cite:
\begin{notebox}
DOI: \url{https://doi.org/10.71674/PH64-N397}
\end{notebox}
\begin{lstlisting}[language=TeX]
@Article{Gao2024,
  author    = {Gao, Stephen and Maurer, Jeremy and
               Information Technology Research Support Solutions},
  title     = {The Mill HPC Cluster},
  year      = {2024},
  doi       = {10.71674/PH64-N397},
  publisher = {Missouri University of Science and Technology},
  url       = {https://scholarsmine.mst.edu/the-mill/1/},
}
\end{lstlisting}
%=============================================================================
\section*{Appendix: The Mill Hardware}
%=============================================================================
\begin{table}[h]
\centering
\begin{tabular}{@{}llll@{}}
\toprule
\textbf{Model} & \textbf{CPU Cores} & \textbf{Memory} & \textbf{Node Count} \\
\midrule
Dell R6525 & 128 & 512 GB & 25 \\
Dell C6525 & 64 & 256 GB & 160 \\
Dell C6420 & 40 & 192 GB & 44 \\
\bottomrule
\end{tabular}
\caption{Compute nodes available on The Mill}
\end{table}
\begin{table}[h]
\centering
\begin{tabular}{@{}lllll@{}}
\toprule
\textbf{Model} & \textbf{CPU Cores} & \textbf{Memory} & \textbf{GPU} & \textbf{GPU Count} \\
\midrule
Dell XE9680 & 112 & 1 TB & H100 SXM5 & 8 \\
Dell C4140 & 40 & 192 GB & V100 SXM2 & 4 \\
Dell R740xd & 40 & 384 GB & V100 PCIe & 2 \\
\bottomrule
\end{tabular}
\caption{GPU nodes available on The Mill}
\end{table}
\end{document}
